\section{Robustness of the neural posteriors}\label{app:robust}

The posteriors in this paper are neural density estimates, which we train by neural posterior estimation
\citep{Papamakarios2016,Greenberg2019} with a masked autoregressive flow \citep{Papamakarios2017}, using the
\textsc{sbi} package \citep{Tejero2020}. A trained network depends on its random seed, which sets the
initial weights and the random choices made during training, and on the number of training simulations.
In this appendix we check that the map-level results at 10 per cent of the sky depend on neither, and we
test the joint posterior in all ten parameters.

\paragraph*{Ensembles.} We train five networks on the same 10\,000 map-level simulations, differing only in
their seed, and combine them into an equal mixture as in equation~(\ref{eq:ensemble}). All ensembles in this
paper, with the numbers of networks given in Tables~\ref{tab:bias} and \ref{tab:coverage}, are built in
this way.

\begin{figure}
\centering
\includegraphics[width=\columnwidth]{figures/ensemble.png}
\caption{Bias $\mathcal{B}$ of the map-level posterior at the true parameters (10 per cent of the sky), in
posterior widths, for $\sigma_8$, $\omc$, $\asz$ and $\alsz$. Light blue circles show five networks trained
on the same 10\,000 simulations with different seeds, and black diamonds their equal mixture,
equation~(\ref{eq:ensemble}). The error bars are the errors of the
mean over 100 map-level mocks, and the dotted line marks zero bias.}
\label{fig:ensemble}
\end{figure}

\begin{table}
\centering
\caption{Robustness of the map-level posterior at 10 per cent of the sky, for five networks trained on
10\,000 simulations with different seeds (smallest and largest value), their ensemble, and one network
trained on 5000 simulations. The top block gives the bias $\mathcal{B}$ in posterior widths, on 100 mocks
at the true parameters, and the bottom block the 90 per cent coverage $\cov(0.9)$ of $\alsz$ on 500 skies
drawn from the prior. The ensemble's coverage in the other parameters is given in Table~\ref{tab:coverage}.}
\label{tab:robust}
\small
\setlength{\tabcolsep}{3pt}
\begin{tabular}{lccc}
\toprule
 & seeds (min, max) & ensemble & 5000 sims \\
\midrule
\multicolumn{4}{l}{\emph{bias $\mathcal{B}$}}\\
$\sigma_8$ & $-0.17$, $+0.02$ & $-0.05\pm0.10$ & $-0.05$ \\
$\omc$     & $+0.02$, $+0.16$ & $+0.10\pm0.07$ & $-0.01$ \\
$\asz$     & $-0.13$, $+0.06$ & $-0.05\pm0.08$ & $+0.02$ \\
$\alsz$    & $+0.01$, $+0.20$ & $+0.08\pm0.08$ & $+0.16$ \\
\midrule
\multicolumn{4}{l}{\emph{coverage $\cov(0.9)$}}\\
$\alsz$    & 0.84, 0.86 & 0.86 & 0.85 \\
\bottomrule
\end{tabular}
\end{table}

The biases of the five networks and of their ensemble (Fig.~\ref{fig:ensemble} and
Table~\ref{tab:robust}) show that the bias of a single network moves by about 0.1 widths from seed to seed.
This scatter is comparable to the error of the mean from 100 mocks. We therefore read the five values for each
parameter as repeated measurements of the same quantity, whose spread is the scatter due to training
alone. In particular, the largest single values, $-0.17$ in $\sigma_8$ and $+0.20$ in $\alsz$, reflect this
scatter rather than a bias of the method. The ensemble (black) is within one error of zero in $\sigma_8$,
$\asz$ and $\alsz$, and gives $+0.10\pm0.07$ in $\omc$.

\paragraph*{Number of simulations.} A single network trained on half the simulations, 5000, has biases of
$-0.05$, $-0.01$, $+0.02$ and $+0.16$ widths in $\sigma_8$, $\omc$, $\asz$ and $\alsz$, respectively. These
values are comparable to the seed-to-seed scatter of about 0.1 widths among the networks trained on 10\,000. In
every case the $\alsz$ intervals are slightly too narrow at 10 per cent of the sky, with 90 per cent
coverage of 0.84--0.86 for the five networks, 0.86 for their ensemble and 0.85 for the network trained on
5000 simulations (Table~\ref{tab:robust}). Since doubling the number of simulations does not change this
coverage, more simulations would probably not remove the deficit. At 40.7 per cent of the sky, by
contrast, the same intervals have correct coverage (0.88, Table~\ref{tab:coverage}).

\paragraph*{Joint coverage: TARP.} Whereas the coverage of equation~(\ref{eq:coverage}) tests one parameter
at a time, the test of accuracy with random points \citep[TARP;][]{Lemos2023} tests the joint posterior in
all ten parameters. We first rescale the parameters to the unit interval over the range of the true values. For each prior draw $n$, we then draw a reference point $\bm\theta_{\mathrm{ref},n}$ uniformly in
that range and compute
\begin{equation}
f_n = \frac{1}{K}\sum_{k=1}^{K}\mathbf 1\!\left[\,\lVert\bm\theta_{n,k}-\bm\theta_{\mathrm{ref},n}\rVert
<\lVert\bm\theta^\star_n-\bm\theta_{\mathrm{ref},n}\rVert\right],
\label{eq:tarp_f}
\end{equation}
the fraction of the $K$ posterior samples $\bm\theta_{n,k}$ that lie closer to the reference point than the
truth $\bm\theta^\star_n$ does. The expected coverage probability at credible level $\gamma$,
$\mathrm{ECP}(\gamma)$, is then the fraction of the $N_\mathrm{draw}$ draws with $f_n<\gamma$, and it
equals $\gamma$ for a posterior with correct coverage. We summarise it by the signed area
\begin{equation}
A_\mathrm{TARP} = \int_{1/2}^{1}\left[\mathrm{ECP}(\gamma)-\gamma\right]\mathrm{d}\gamma,
\label{eq:tarp_a}
\end{equation}
which is 0 for a posterior with correct coverage, negative for one that is too narrow and positive for one
that is too wide. Using the implementation in \textsc{sbi} on the 500 prior skies at 10 per cent of the
sky, we find $A_\mathrm{TARP}$ between $-0.009$ and $-0.029$ for the five map-level networks. The joint
posterior is therefore slightly too narrow, consistent with the 90 per cent coverage of 0.84--0.86 in
$\alsz$ for the same networks.
